\begin{thebibliography}{99}

\bibitem{huang1998}
N. E. Huang \emph{et al.}, ``The empirical mode decomposition and the Hilbert spectrum for nonlinear and non-stationary time series analysis,'' \emph{Proc. R. Soc. Lond. A}, vol.~454, no.~1971, pp. 903--995, 1998.

\bibitem{schalk2004}
G. Schalk, D. J. McFarland, T. Hinterberger, N. Birbaumer, and J. R. Wolpaw, ``BCI2000: A general-purpose brain-computer interface (BCI) system,'' \emph{IEEE Trans. Biomed. Eng.}, vol.~51, no.~6, pp. 1034--1043, 2004.

\bibitem{eegmmidb}
G. Schalk, ``EEG Motor Movement/Imagery Dataset,'' version 1.0.0, PhysioNet, 2009, doi: 10.13026/C28G6P.

\bibitem{goldberger2000}
A. L. Goldberger \emph{et al.}, ``PhysioBank, PhysioToolkit, and PhysioNet: Components of a new research resource for complex physiologic signals,'' \emph{Circulation}, vol.~101, no.~23, pp. e215--e220, 2000.

\bibitem{pollard2026}
T. Pollard, B. E. Moody, L. Lehman, B. Gow, C. Fernandes, C. Xie, A. Johnson, R. G. Mark, and T. Heldt, ``PhysioNet as a global platform for biomedical research,'' \emph{Nature Health}, vol.~1, no.~8, pp. 792--795, 2026.

\bibitem{zyma2019}
I. Zyma \emph{et al.}, ``Electroencephalograms during mental arithmetic task performance,'' \emph{Data}, vol.~4, no.~1, p.~14, 2019.

\bibitem{eegmat}
I. Zyma, I. Seleznov, A. Popov, M. Chernykh, and O. Shpenkov, ``EEG During Mental Arithmetic Tasks,'' version 1.0.0, PhysioNet, 2018, doi: 10.13026/C2JQ1P.

\bibitem{brunner2008}
C. Brunner, R. Leeb, G. R. M{\"u}ller-Putz, A. Schl{\"o}gl, and G. Pfurtscheller, ``BCI Competition 2008 -- Graz data set A,'' Graz University of Technology, 2008.

\bibitem{tangermann2012}
M. Tangermann \emph{et al.}, ``Review of the BCI Competition IV,'' \emph{Front. Neurosci.}, vol.~6, p.~55, 2012.

\bibitem{laszuk2017}
D. Laszuk, ``Python implementation of Empirical Mode Decomposition algorithm,'' GitHub, 2017, doi: 10.5281/zenodo.5459184.

\bibitem{welch1967}
P. D. Welch, ``The use of fast Fourier transform for the estimation of power spectra: A method based on time averaging over short, modified periodograms,'' \emph{IEEE Trans. Audio Electroacoust.}, vol.~15, no.~2, pp. 70--73, 1967.

\bibitem{ledoit2004}
O. Ledoit and M. Wolf, ``A well-conditioned estimator for large-dimensional covariance matrices,'' \emph{J. Multivariate Anal.}, vol.~88, no.~2, pp. 365--411, 2004.

\bibitem{itakura1968}
F. Itakura and S. Saito, ``Analysis synthesis telephony based on the maximum likelihood method,'' in \emph{Proc. 6th Int. Congr. Acoust.}, Tokyo, Japan, 1968, pp. C17--C20.

\bibitem{pedregosa2011}
F. Pedregosa \emph{et al.}, ``Scikit-learn: Machine learning in Python,'' \emph{J. Mach. Learn. Res.}, vol.~12, pp. 2825--2830, 2011.

\bibitem{efron1993}
B. Efron and R. J. Tibshirani, \emph{An Introduction to the Bootstrap}. New York, NY, USA: Chapman \& Hall, 1993.

\bibitem{benjamini1995}
Y. Benjamini and Y. Hochberg, ``Controlling the false discovery rate: A practical and powerful approach to multiple testing,'' \emph{J. R. Statist. Soc. B}, vol.~57, no.~1, pp. 289--300, 1995.

\bibitem{wilcoxon1945}
F. Wilcoxon, ``Individual comparisons by ranking methods,'' \emph{Biometrics Bull.}, vol.~1, no.~6, pp. 80--83, 1945.

\bibitem{holm1979}
S. Holm, ``A simple sequentially rejective multiple test procedure,'' \emph{Scand. J. Statist.}, vol.~6, no.~2, pp. 65--70, 1979.

\bibitem{pfurtscheller1999}
G. Pfurtscheller and F. H. Lopes da Silva, ``Event-related EEG/MEG synchronization and desynchronization: Basic principles,'' \emph{Clin. Neurophysiol.}, vol.~110, no.~11, pp. 1842--1857, 1999.

\bibitem{ang2012}
K. K. Ang, Z. Y. Chin, C. Wang, C. Guan, and H. Zhang, ``Filter bank common spatial pattern algorithm on BCI Competition IV datasets 2a and 2b,'' \emph{Front. Neurosci.}, vol.~6, p.~39, 2012.

\end{thebibliography}
